% ==============================================================================
% Secção 1: Introdução
% ==============================================================================

\subsection{Título do projeto (em português)}
\begin{campotexto}
Otimização do Carro do Futuro (Beyond Wheels)
\end{campotexto}

\subsection{Título do projeto (em inglês)}
\begin{campotexto}
Future Cars Optimization (Beyond Wheels)
\end{campotexto}

\subsection{Acrónimo}
\begin{campotexto}
BW
\end{campotexto}

\subsection{Palavras-chave (em português)}
\begin{campotexto}
Bateria de Lítio, Bateria de Sódio, Suspensão Adaptativa, Regeneração de Energia
\end{campotexto}

\subsection{Palavras-chave (em inglês)}
\begin{campotexto}
Lithium Battery, Sodium Battery, Adaptive Suspension, Energy Regeneration
\end{campotexto}

\subsection{Áreas científicas principais}
\begin{campotexto}
Engenharia Mecânica, Eletrónica
\end{campotexto}

\subsection{Duração em meses}
\begin{campotexto}
36
\end{campotexto}

\subsection{Identificação de questões éticas}
\begin{campotexto}
As questões éticas do BeyondWheels resultam sobretudo do impacto ambiental e
social das tecnologias estudadas e da segurança de quem utiliza os veículos.
No armazenamento de energia, a transição das baterias de lítio para baterias de
sódio e a sua reciclagem no fim de vida respondem a preocupações éticas e
ambientais associadas à extração de lítio, ao descarte de baterias e ao risco de
incêndio. É assim necessário garantir que a cadeia de valor dos materiais e o
tratamento das baterias usadas sejam sustentáveis e responsáveis.

A suspensão adaptativa e preditiva, baseada em sensores e Inteligência
Artificial, atua diretamente sobre a estabilidade e o conforto do veículo.
Decisões incorretas ou falhas do sistema podem, por isso, afetar a segurança de
passageiros e de outros utilizadores da estrada, o que exige validação rigorosa,
modos de funcionamento seguro em caso de falha e conformidade com a
regulamentação e a certificação de segurança europeias. Por recolher e
processar dados em tempo real sobre o veículo e a via, o sistema levanta ainda
questões de privacidade e de proteção de dados, devendo respeitar o RGPD e ser
transparente quanto aos dados recolhidos, à sua finalidade e às decisões
tomadas pelo modelo.
\end{campotexto}

\subsection{Caracterização dos proponentes}
\begin{campotexto}
Nesta secção apresenta-se a equipa proponente do projeto BeyondWheels, com as
suas competências, funções e contributos para o desenvolvimento da proposta. O
projeto é interdisciplinar, ao integrar mecânica, eletrónica, materiais e
sustentabilidade e modelos de Inteligência Artificial. A
equipa foi, por isso, organizada por funções, cada uma associada a um conjunto
de atividades distinto: a definição do método das suspensões adaptativas, a
regeneração de energia, a transição das baterias de lítio para sódio, a
autonomia e reciclagem das baterias, o orçamento e a comunicação do projeto
através da página web.

O projeto desenvolve-se em consórcio com dois outros grupos, que colaboram de
forma localizada em atividades específicas, mantendo o Grupo 3 a coordenação e
a responsabilidade pela proposta:
\begin{itemize}
  \item \textbf{Grupo 2}, cujo tema são veículos elétricos autónomos que
  circulam numa via magnética dedicada em autoestrada e passam a rodas em
  estrada convencional. É nesta transição que entram em funcionamento as
  baterias e as suspensões, tornando este grupo um parceiro de aplicação e
  validação das soluções propostas.
  \item \textbf{Grupo 4}, cujo tema são suspensões adaptativas e preditivas, adaptando-se ao perfil da estrada,
  que contribui para o desenvolvimento da tecnologia de suspensão inteligente
  prevista no projeto.
\end{itemize}

Os contributos dos parceiros são complementares e não substituem o trabalho
desenvolvido pelo Grupo 3. De seguida descreve-se o envolvimento de cada
elemento do grupo proponente, com a identificação das suas competências e
funções no projeto.
\end{campotexto}

% --- Macro auxiliar para tabelas de elementos internos ---
\newcommand{\tabelaElemento}[5]{%
  \begin{table}[H]
  \centering\small
  \begin{tabularx}{\textwidth}{L{0.32\textwidth} Y}
    \toprule
    \rowcolor{uatablehead}
    \multicolumn{2}{l}{\sffamily\bfseries\color{uadark} #1} \\
    \midrule
    \sffamily\bfseries Identificação & #2 \\ \arrayrulecolor{uaborder}\hline
    \sffamily\bfseries Competências & #3 \\ \hline
    \sffamily\bfseries Funções no projeto & #4 \\ \hline
    \sffamily\bfseries Envolvimento no desenvolvimento da proposta (tópicos e percentagem global) & #5 \\
    \arrayrulecolor{black}\bottomrule
  \end{tabularx}
  \end{table}
}

% --- Macro auxiliar para tabelas de entidades externas ---
\newcommand{\tabelaParceiro}[7]{%
  \begin{table}[H]
  \centering\small
  \begin{tabularx}{\textwidth}{L{0.32\textwidth} Y}
    \toprule
    \rowcolor{uatablehead}
    \multicolumn{2}{l}{\sffamily\bfseries\color{uadark} #1} \\
    \midrule
    \sffamily\bfseries Identificação & #2 \\ \arrayrulecolor{uaborder}\hline
    \sffamily\bfseries Representante & #3 \\ \hline
    \sffamily\bfseries Competências & #4 \\ \hline
    \sffamily\bfseries Funções no projeto & #5 \\ \hline
    \sffamily\bfseries Entidade / setor & #6 \\ \hline
    \sffamily\bfseries Elementos que constituem o proponente e respetiva função (facultativo) & #7 \\
    \arrayrulecolor{black}\bottomrule
  \end{tabularx}
  \end{table}
}

\tabelaElemento{Elemento 1 -- Entidade proponente \#1 (grupo 3)}
{{Diana Baldaia}}
{{Gestão de equipas, Mecânica e Modelos de IA}}
{{Chefe de projeto; proponente da ideia das suspensões adaptativas inteligentes e definição do respetivo método}}
{{Conceito e método das suspensões adaptativas, coordenação da equipa e redação da proposta, 30\%}}

\tabelaElemento{Elemento 2 -- Entidade proponente \#1 (grupo 3)}
{{Leandro Simões}}
{{Trabalho em equipa}}
{{Apoio pontual em tarefas atribuídas ao longo do projeto}}
{{Apoio em tarefas específicas da proposta, 8\%}}

\tabelaElemento{Elemento 3 -- Entidade proponente \#1 (grupo 3)}
{{Marco Pereira}}
{{Eletrónica, Competências Sociais}}
{{Proponente da ideia da regeneração de energia nas suspensões e definição do respetivo processo; gestão do consórcio, divisão de tarefas e acompanhamento do trabalho da equipa}}
{{Regeneração de energia, organização e distribuição de tarefas, 20\%}}

\tabelaElemento{Elemento 4 -- Entidade proponente \#1 (grupo 3)}
{{Mariana Almeida}}
{{Gestão orçamental}}
{{Responsável pela transição de baterias de Lítio para baterias de Sódio e pela elaboração dos orçamentos gerais}}
{{Transição para baterias de Sódio e orçamento do projeto, 12\%}}

\tabelaElemento{Elemento 5 -- Entidade proponente \#1 (grupo 3)}
{{Rafael Alegria}}
{{Design e desenvolvimento web}}
{{Web Designer; contributo para a ideia das baterias e responsável pela página web do projeto}}
{{Conceção e gestão da página web, 12\%}}

\tabelaElemento{Elemento 6 -- Entidade proponente \#1 (grupo 3)}
{{Simão Pires}}
{{Gestão de projeto, Materiais e Sustentabilidade}}
{{Autonomia das baterias, reciclagem das baterias em fim de vida, sustentabilidade e planeamento do projeto}}
{{Baterias, materiais, sustentabilidade e planeamento, 18\%}}

\tabelaParceiro{Entidade externa \#1 (Grupo 2)}
  {MagLink}
  {Matilde Romano}
  {Otimização no mercado dos transportes}
  {Implementação das melhorias sugeridas pela equipa de investigação e respetivo feedback}
  {Setor dos transportes e desenvolvimento de novas tecnologias}
  {Diogo Pedrosa, Gabriel Ribeiro, Guilherme Lopes, Lucas Fonseca e Matilde Romano}

\tabelaParceiro{Entidade externa \#2 (Grupo 4)}
  {Suspensão Ativa Preditiva}
  {Dinis Mano}
  {Conhecimentos em engenharia mecânica e sistemas inteligentes}
  {Apoio ao desenvolvimento das tecnologias de suspensão}
  {Investigação aplicada em engenharia mecânica e sistemas inteligentes}
  {André Pinto, Dinis Mano, João Martins, Luis Guimarães e Rodrigo Osório}
